\documentclass[10pt]{article}

\usepackage[utf8]{inputenc}

\usepackage[T1]{fontenc}

\usepackage{amsmath}

\usepackage{amsfonts}

\usepackage{amssymb}

\usepackage[version=4]{mhchem}

\usepackage{stmaryrd}

\usepackage{bbold}



\begin{document}

Обобще н ны й ме тод годографа - метод, применимый к системам гидродинамич. типа с диагональной матрицей



\[

u_{t}^{i}=v_{i}\left(u^{1}, \ldots, u^{n}\right) u_{x}^{i}, u^{i}=u^{i}(x, t), i=1, \ldots, n,

\]



и позволяющий свести систему (2) к линейной системе уравнений



\[

\partial_{i} w_{k}=\frac{\partial_{i} v_{k}}{\left(v_{i}-v_{k}\right)}\left(w_{i}-w_{k}\right), \quad \partial_{i}=\frac{\partial}{\partial u^{i}}, \quad i \neq k,

\]



относительно коэффициентов $w_{i}\left(u^{1}, \ldots, u^{n}\right)$ коммутирующих с (2) потоков (то есть симметрий (2)) того же вида при выполнении условий «полугамильтоновости» (см. [1]):



\[

\partial_{i}=\left(\frac{\partial_{j} v_{k}}{v_{j}-v_{k}}\right)=\partial_{j}\left(\frac{\partial_{i} v_{k}}{v_{i}-v_{k}}\right), \quad i \neq j \neq k .

\]



Связь решений (2) и (3) дает формула обобщенного преобразования годографа



\[

w_{i}(u)=v_{i}(u) t+x

\]



Jum.: [1] М и з с Р., Математическая теория течений сжимаемой жидкости, пер. с англ., М., 1961; [2] Куран т Р., Уравнения с частными производными, пер. с англ., М., 1964; [3] Дубров и н Б. А., Н о в и ко в С. П., "Успехи математич. наук», 1989, т. 44, $\begin{array}{ll}\text { в. 6, с. 29-98. } & \text { С. П. Царев. }\end{array}$ ГОДОГРАФА ПРЕОБРАЗОВАНИЕ - см. Годографа меmod.



ГОЛДЕНА - TOMПСОНА НЕРАВЕНСТВО - неравенство для положительных самосопряженных операторов, определяемое следующим образом. Пусть $H$ - сепарабельное гильбертово пространство, а $\boldsymbol{E}_{p}(H)$ - банахово пространство компактных операторов с конечной $\|\cdot\|_{p}$-нормой:



\[

\|A\|_{p}=\left(\sum_{n=1}^{\infty} \lambda_{n}^{p}(A)\right)^{1 / p}

\]



где $\lambda_{n}(A)$ - собственные значения оператора $|A|=\left(A^{*} A\right)^{1 / 2}$. Если $A$ и $B$ - положительные самосопряженные операторы, оператор $A+B$ в существенном сопряжен на их общей области определения $D(A) \cap D(B)$ и, по крайней мере, один из них порождает гиббсовскую полугруппу, тогда



\[

\exp \left\{-(A+B)^{\sim}\right\} \in \boldsymbol{S}_{1}(H)

\]



и имеет место неравенство Голдена - Томпсон а (см. [1], [2]):



$\operatorname{tr} \exp \left\{-(A+B)^{\sim}\right\} \leqslant \operatorname{tr}[\exp (-A) \exp (-B)]$,



что эквивалентно



\[

\begin{gathered}

\left\|\exp \left\{-(A+B)^{\sim}\right\}\right\|_{1} \leqslant \\

\leqslant\|\exp (-A / 2) \exp (-B) \exp (-A / 2)\|_{1} .

\end{gathered}

\]



Обобщение Г.- Т. н. для $p>1$ :



\[

\begin{gathered}

\left\|\exp \left\{-(A+B)^{\sim}\right\}\right\|_{p} \leqslant \\

\leqslant \exp (-A / 2) \exp (-B) \exp (-A / 2) \|_{p}

\end{gathered}

\]



содержится в работе [3], см. также [4]. Для $p=\infty$ это неравенство известно как лемма Сигала (см. [4]). С помощью одного неравенства Вейля Г.- Т. Н. можно обобшить на случай произвольной функции $\varphi$ такой, что отображение $t \rightarrow \varphi\left(e^{t}\right)$ выпукло (см. [4]):



\[

\begin{gathered}

\operatorname{tr} \varphi\left[\exp \left\{-(A+B)^{\sim}\right\}\right] \leqslant \\

\leqslant \operatorname{tr} \varphi[\exp (-A / 2) \exp (-B) \exp (-A / 2)] .

\end{gathered}

\]



Jum.: [1] G old e n S., «Phys. Rev.», 1965, v. 137B, № 4, p. 1127-28; [2] Thomps on C., «J. Math. Phys.», 1965, v. 6, № 11, p. 181-213; [3] Ruska i M. B., “Comm. Math. Phys.», 1972, v. 26, № 4, p. 280-89; [4] Рид М., С аймон Б., Методы современной математической физики, т. 4: Анализ операторов, пер. с англ., М., 1982.



B. A. Загребнов. ГОЛДСТОУНА ТЕOPEMА: спонтанное нарушение непрерывной глобальной симметрии приводит к появлению безмассовой бесспиновой частицы. Установлена в [1]; эта безмассовая частица называется голдсто у о вским 6 о30 ном. Спонтанное нарушение симметрии, конечно, остается симметрией системы. Но эта симметрия проявляется не в инвариантности основного состояния относитель- но преобразований этой симметрии, а в появлении го л д с т о но в с ко й моды (существовании голдстоуновского бозона).



Примерами голдстоуновских возбуждений в нерелятивистской теории многих тел являются спиновые волны в ферромагнетике, фононы в кристаллич. твердых телах и жидком гелии.



Jum.: [1] G oldsto ne J., "Nuovo Cim.», 1961, v. 19, p. 154.



ГолДСТОУНОВСКАЯ МОДА - см. Голдстоуна теорема, Спонтанное нарушение симметрии.



ГолДСтоУНОвСКМЙ БОзОН - см. Голдстоуна теорема. ГОЛОВАСТИК - см. Фейнмана диаграмма.



ГологРАФия - возможность записи и воспроизведения полей. Г. обнаружена Д. Габором [1] для записи в плоской светочувствительной среде. Бурное развитие Г. последовало после того, как в 1962-64 были применены лазер и внеосевая схема записи. Двумерная Г. оказалась лишь частным случаем трехмерной, обладающей более полным комплексом отображаемых свойств (см. [2]). В многочисленных разделах Г. изучаются типы и виды голограмм, формирование изображений, аппаратура, методы, теория спеклона и т. д. (см. [3]-[5]).



Принципы Г. подробно изложены в физич. литературе, однако строгое описание явлений, происходящих при плоской Г., сталкивается с трудностями, связанными с отсутствием математич. постановки задачи рассеяния волн на полупрозрачном экране, что имеет прямое отношение к проблеме «черного тела» (см. [6]). Поставить задачу введением каких-либо условий на экране не удается. Эту трудность можно преодолеть с помощью введения подходяших условий поглощения на экране (см. [7]). С их помощью определяется характеристика пропускания голограммы и рассчитывается рассеянная волна. При этом применяется теория коротковолновых асимптотик и фронтов осцилляций по параметру $\lambda=c / \omega, \lambda \rightarrow 0$ (физически - длина волны $\lambda$ много меньше характерного размера $L$ ).



Пусть оператор $\Gamma^{S} \equiv \Gamma_{\lambda}^{S}$ определен формулой



\[

\Gamma^{S}(\psi)=\int_{S} G(x, y) M\left(n_{y}\right) \psi(y) d \sigma(y) .

\]



Здесь $n_{y}-$ внешняя нормаль к поверхности $S, d \sigma-$ элемент площади поверхности, $G=G_{\lambda}$ обозначает удовлетворяющее условиям излучения левое фундаментальное решение для системы Максвелла,



\[

M(p)=\left\|\begin{array}{cc}

0 & -\Lambda(p) \\

\Lambda(p) & 0

\end{array}\right\|

\]



где $\Lambda(p)=[p \times \ldots]$ - матрица векторного умножения на $p \in \mathbb{R}^{3}$.



На экран падает волна



\[

\Phi(t)=\int \varphi_{\omega} \exp (-i \omega t) d \omega .

\]



У с л о в и е А. Волна $\Phi_{+}$, прошедшая экран, задается формулой



\[

\Phi_{+}=\int \Gamma_{c / \omega}^{S}\left(f_{\omega} \varphi_{\omega}\right) \exp (-i \omega t) d \omega,

\]



где $f_{\omega}-$ скалярная функция на экране (характеристика пропускания на частоте $\omega)$. Область, где $f_{\omega}=1$, прозрачна для волн с частотой $\omega$, а область, где $f_{\omega}=0$, абсолютно поглощает.



Пусть экран $S$ получен как позитивный отпечаток (то есть на $S$ помещается светочувствительный материал и освешается световой волной $\Phi$; по негативному отпечатку изготовляется позитивный).



У с л о в и е Б. Характеристика пропускания полученного экрана имеет вид:



\[

f_{\omega}=\rho^{1}(y) F_{\omega}\left(\int \rho^{2}(t) I(y, t) d t\right),

\]



где $I=\Phi \Phi^{*}-$ интенсивность светового импульса, функции $\rho^{1}, \rho^{2} \in C_{0}^{\infty}$ описывают переходные процессы. Предполагается, что существует интервал, где функция $F_{\omega}$ линейна.





\end{document}